\documentclass[10pt,a4paper]{article}

\usepackage[T1]{fontenc}
\usepackage[utf8]{inputenc}
\usepackage{lmodern}
\usepackage{amsmath,amssymb,amsthm}
\usepackage{booktabs,array}
\usepackage{geometry}
\usepackage{xcolor}
\usepackage{enumitem}
\usepackage{microtype}
\usepackage{hyperref}
\usepackage{fancyhdr}

\geometry{margin=1in}
\hypersetup{
  colorlinks=true,
  linkcolor=blue!55!black,
  citecolor=blue!55!black,
  urlcolor=blue!55!black,
  pdftitle={A Quarter Turn Is All It Takes: A Zero-Query Forgery against Yagisawa's Quaternion Signature},
  pdfauthor={ePrint Signature Audit}
}
\setlist{nosep,leftmargin=*}
\setlength{\parindent}{1em}
\setlength{\parskip}{0.25em}
\pagestyle{fancy}
\fancyhf{}
\fancyhead[L]{A Quarter Turn Is All It Takes}
\fancyhead[R]{9 October 2026}
\fancyfoot[C]{\thepage}
\renewcommand{\headrulewidth}{0.3pt}

\newtheorem{theorem}{Theorem}
\newtheorem{lemma}{Lemma}
\newcommand{\F}{\mathbb{F}}
\newcommand{\Hq}{\mathbb{H}_{q}}
\newcommand{\vect}[1]{\boldsymbol{#1}}

\title{\bfseries A Quarter Turn Is All It Takes: A Zero-Query Forgery against Yagisawa's Quaternion Signature}
\author{ePrint Signature Audit\\
\small AI-assisted cryptanalysis report}
\date{9 October 2026}

\begin{document}
\maketitle

\noindent\fcolorbox{black!50}{black!4}{\parbox{0.94\textwidth}{
\textbf{AI disclosure and review status.}
This paper documents this audit's AI-assisted reconstruction of the attack.
OpenAI Codex performed exact-version pinning, attack-history searching, scheme
reconstruction, cryptanalysis, proof development, implementation, experiment
generation, and technical writing.  Specialist AI agents were used for
technical, adversarial, reproducibility, and publication review.  Independent
AI reproduction is identified as AI reproduction and is not described as
human verification.  OpenAI Codex generated and selected the title and
intentionally chose ``A Quarter Turn Is All It Takes'' as wordplay on the
quaternion/order-four mechanism: multiplication by the basis element $i$ is a
quarter turn, and its four-point power orbit defeats the verifier.  The
technical, adversarial, independent-reproduction, final-history, and
publication-review gates passed; this result is classified
\texttt{VERIFIED\_ATTACK}.  No claim of independent human verification is
made.}}

\begin{abstract}
We give a public-key-only, zero-query forgery against Masahiro Yagisawa's
\emph{A Digital Signature Using Multivariate Functions on Quaternion Ring},
IACR ePrint 2010/352, final revision \texttt{20100627:124304}.  For the paper's
claimed row $d=2$, $r=3$, and $m=448$, the public quaternion-valued polynomial
has homogeneous degree $s=13$ and 2,240 field coefficients.  From those public
coefficients alone, the adversary chooses $R=i$ and the fresh message $E=2k$,
so $g=2$, and transmits
\[
  T^*(X)=F(iX)+(x_2^{13},0,0,0).
\]
The composition $F(iX)$ is only a signed permutation of public coefficients.
The added monomial vanishes on every verifier challenge point
$R^{p+1}\in\{1,i,-1,-i\}$, while it is nonzero at the verifier's separate
inequality point $RE=-2j$.  The ordinary verifier therefore accepts for every
integer challenge $p$.

At $q=1{,}048{,}609$, a full-row reference implementation generated a
legitimate key, forged in 0.001789767 seconds, and completed key generation and
all checks in 112.834013369 seconds.  External timing was 1:53.69 with
19,968~KiB peak resident memory.  Honest and forged signatures on the same
message each passed 80/80 sampled challenges, all four challenge residue
classes passed exactly, and 7/7 negative controls rejected.  The forger made
no signing query.  We located no previous public attack against this exact
construction in the searches performed as of 2026-10-09.  We make no novelty
claim for the general vanishing-polynomial or small-evaluation-set principle.
\end{abstract}

\section{Introduction}

Yagisawa proposes a signature construction over a quaternion ring over a
finite field~\cite{yagisawa}.  The public key is a four-coordinate homogeneous
polynomial $F$.  A signature carries a second four-coordinate polynomial $T$,
a quaternion $R$, and a message quaternion $E$.  Verification does not check
that $T$ was derived from the secret key.  It instead evaluates $F$ and $T$ at
one message-dependent point and one random power of $R$.

The paper already observes that an adversary can publicly form the shifted
polynomial $F(R^{g-1}X)$.  Its preliminary inequality is intended to reject
that construction.  The check has only one point, however.  We add a
homogeneous polynomial that is zero on every random-challenge point but
nonzero at the preliminary point.  Choosing the order-four quaternion $i$
makes both the construction and its proof especially simple.

Our concrete result has the following properties:
\begin{itemize}
  \item it attacks the final ePrint revision and the paper's complete claimed
  parameter row, without weakening $q,d,r$, or $m$;
  \item it uses the ordinary two-stage verifier and a normally signable
  message, confirmed by an honest signature on that same message;
  \item it transforms public coefficients only and makes zero signing queries;
  \item its randomized verification equality holds for every challenge $p$;
  \item it is a direct fresh-message forgery, rather than key recovery; and
  \item its online work is a permutation and sign change of 2,240 field
  coefficients followed by one field addition.
\end{itemize}

\paragraph{Attack-history statement.}
We located no previous public attack against this exact Yagisawa construction
in the searches performed as of 2026-10-09.  This is a dated statement about
the searches conducted, not an absolute priority claim.  Published attacks on
other quaternion signatures concern different constructions.

\paragraph{Technique attribution.}
The broad fact that a nonzero polynomial can evade evaluation on a small or
structured test set underlies classical polynomial identity testing, including
the work of Zippel~\cite{zippel} and Schwartz~\cite{schwartz}.  Our attack is
an elementary, exact vanishing-polynomial construction on the verifier's
four-point orbit.  We claim an application to this target, not novelty for the
technique.

\section{The target construction}

\subsection{Quaternion convention}

Let $q$ be an odd prime and write an element of $\Hq$ as
$A=(a_0,a_1,a_2,a_3)=a_0+a_1i+a_2j+a_3k$.  The paper defines the product of
$A$ and $B=(b_0,b_1,b_2,b_3)$ by
\[
AB=\left(\begin{array}{l}
a_0b_0-a_1b_1-a_2b_2-a_3b_3,\\
a_0b_1+a_1b_0+a_2b_3-a_3b_2,\\
a_0b_2-a_1b_3+a_2b_0+a_3b_1,\\
a_0b_3+a_1b_2-a_2b_1+a_3b_0
\end{array}\right) \pmod q.
\]
Thus $i^2=j^2=k^2=-1$, $ij=k$, $ik=-j$, and $ki=j$.

\subsection{Keys, signatures, and verification}

For secret $k_1,\ldots,k_m\in\F_q$ and
$A_1,\ldots,A_m\in\Hq$, the public polynomial is
\begin{equation}
 F(X)=\sum_{h=1}^{m} k_h
       \prod_{\ell=0}^{d}\left(A_h^{r^\ell}X^{r^\ell}\right),
 \label{eq:F}
\end{equation}
where the noncommutative product is taken in increasing $\ell$.  If
\[
 s=1+r+\cdots+r^d,
\]
then each coordinate of $F$ is published as a homogeneous commutative
polynomial
\[
 f_u(x_0,x_1,x_2,x_3)
 =\sum_{e_0+e_1+e_2+e_3=s}
 f_{u,e_0,e_1,e_2,e_3}
 x_0^{e_0}x_1^{e_1}x_2^{e_2}x_3^{e_3}.
\]

For a message $E=(E_0,E_1,E_2,E_3)$, set
$g=E_0+E_1+E_2+E_3$.  An honest signer selects $R$ and forms
\begin{equation}
 T(X)=\sum_{h=1}^{m} k_h
       \prod_{\ell=0}^{d}
       \left(A_h^{r^\ell}R^{(g-1)r^\ell}X^{r^\ell}\right).
 \label{eq:T}
\end{equation}
The target paper writes the transmitted signed object as the coefficients of
$T$, together with $R$ and $E$.

The ordinary verifier first requires
\begin{equation}
 F(R^gE)\ne T(RE),
 \label{eq:inequality}
\end{equation}
then samples an integer $p$ and accepts if
\begin{equation}
 F(R^{g+p})=T(R^{p+1}).
 \label{eq:challenge}
\end{equation}
The recommended row is
\[
  d=2,\qquad r=3,\qquad m=448,
\]
so $s=1+3+9=13$.  Each coordinate has
$\binom{s+3}{3}=560$ coefficients and the complete public key has
$4\binom{16}{3}=2{,}240$ coefficients.  The paper requires an odd $q$ larger
than 20 bits.

\paragraph{Claimed security.}
For this row, Yagisawa models recovery of the hidden $k_h$ and $A_h$ values as
2,240 degree-14 multivariate equations in 2,240 secret variables.  Using a
degree of regularity 15 and matrix exponent 2.39, the paper estimates the
Gr\"obner-basis work at about $2^{302}$, compares it with a $2^{80}$ security
target, and claims resistance to Gr\"obner-basis and differential attacks.
It also requires $q$ to exceed 20 bits so that the value space for the
verification evaluations exceeds $2^{80}$.  These are the target paper's
claims, not estimates validated here.  The forgery below bypasses recovery of
all 2,240 secret variables and therefore does not test those estimates.

\section{The zero-query forgery}

Set
\[
 R=i=(0,1,0,0),\qquad E=2k=(0,0,0,2),\qquad g=2.
\]
These values do not commute because
\[
 i(2k)=-2j,\qquad (2k)i=2j.
\]
For $X=(x_0,x_1,x_2,x_3)$, left multiplication by $i$ is the linear map
\begin{equation}
 iX=(-x_1,x_0,-x_3,x_2).
 \label{eq:left-i}
\end{equation}
The attacker computes from the public key
\begin{equation}
 T^*(X)=F(iX)+D(X),
 \qquad D(X)=(x_2^{13},0,0,0).
 \label{eq:forgery}
\end{equation}
Both summands are homogeneous of degree 13, so $T^*$ has exactly the required
signature format.

\subsection{Public coefficient transformation}

Equation~\eqref{eq:left-i} maps a public monomial by
\begin{align*}
 x_0^{e_0}x_1^{e_1}x_2^{e_2}x_3^{e_3}
 &\longmapsto
 (-x_1)^{e_0}x_0^{e_1}(-x_3)^{e_2}x_2^{e_3}\\
 &=(-1)^{e_0+e_2}
 x_0^{e_1}x_1^{e_0}x_2^{e_3}x_3^{e_2}.
\end{align*}
Thus every exponent tuple
\[
 (e_0,e_1,e_2,e_3)
 \longmapsto (e_1,e_0,e_3,e_2)
\]
with sign $(-1)^{e_0+e_2}$.  Constructing $F(iX)$ is a signed permutation of
the public coefficient array.  The attacker then adds one modulo $q$ to output
coordinate zero at exponent tuple $(0,0,13,0)$.  No secret quantity or oracle
response appears in this computation.

\subsection{Proof of acceptance}

\begin{theorem}
For every odd prime $q$ and every public polynomial $F$ in the claimed
degree-13 format, the signed object defined in~\eqref{eq:forgery} is accepted
for every integer verifier challenge $p$.
\end{theorem}

\begin{proof}
Because $q$ is odd, $i^2=-1\ne1$ and $i^4=1$, so $R=i$ has exact order four.
For every integer $p$,
\[
 R^{p+1}\in\{1,i,-1,-i\}.
\]
All four points have zero $x_2$ coordinate, hence
$D(R^{p+1})=0$.  Since $g=2$ and $i$ commutes with its powers,
\begin{align*}
 T^*(R^{p+1})
 &=F(iR^{p+1})+D(R^{p+1})\\
 &=F(R^{p+2})\\
 &=F(R^{g+p}).
\end{align*}
The randomized equality~\eqref{eq:challenge} therefore holds for every $p$.

For the preliminary inequality,
\[
 RE=i(2k)=-2j,
 \qquad i(RE)=-2k=R^2E.
\]
Consequently the shifted part agrees with the left side,
\[
 F(i(RE))=F(R^2E),
\]
but
\[
 D(RE)=((-2)^{13},0,0,0)\ne0
\]
over every odd prime field.  Hence
$T^*(RE)\ne F(R^2E)$, so~\eqref{eq:inequality} also passes.
\end{proof}

\paragraph{Freshness and attack model.}
The attacker chooses $E=2k$ and produces $(T^*,R,E)$ from the public
coefficient representation of $F$.  It requests no signature, making $E$
fresh by construction.  Verification uses the two equations printed in the
paper, with no attack-specific acceptance rule.

\section{Message-domain conditions}

The paper asks that $R$ not commute with any hidden $A_h$, and that $E$ not
commute with $R$ or any $A_h$.  The public verifier cannot evaluate the
conditions involving the hidden $A_h$.  They nevertheless hold exactly for
the reproduced key.

For completeness, let $A=(a_0,a_1,a_2,a_3)$.  Direct multiplication shows
\[
 Ai=iA \iff a_2=a_3=0,
 \qquad
 Ak=kA \iff a_1=a_2=0.
\]
The centralizers of $i$ and $k$ therefore each contain $q^2$ elements and
intersect in the $q$ scalar quaternions.  Multiplication by the nonzero scalar
2 does not change the centralizer of $k$.  For a uniform $A$,
\[
 \Pr[A\text{ commutes with }i\text{ or }2k]
 =\frac{2}{q^2}-\frac{1}{q^3}.
\]
For $m$ independently sampled constants, the exact failure probability is
\[
 1-\left(1-\frac{2}{q^2}+\frac{1}{q^3}\right)^m,
\]
and the union bound is
$m(2/q^2-1/q^3)$.  At $q=1{,}048{,}609$ and $m=448$, this is less than
$8.15\times10^{-10}$.  The attack is therefore overwhelmingly applicable
under the natural uniform-key interpretation, and a test instance can audit
the hidden conditions exactly.

\section{Full-parameter reproduction}

The reference implementation uses only Python's standard library.  It first
generates $m=448$ secret terms and expands Equation~\eqref{eq:F} into all
2,240 public coefficients.  Secret-factorized and public-coefficient
evaluations are cross-checked before the secret state is excluded from the
forger interface.  The forger receives only $F$, returns $T^*,R,E$, and
finishes before the honest same-message control is generated.

The retained primary execution used Python~3.12.3 and its standard library on
an AMD EPYC Processor, with one worker.  The independent execution used the
same language/runtime class and worker count, but a separate polynomial
implementation, deterministic context, and conforming prime.

The same run then constructs Equation~\eqref{eq:T} honestly for $E=2k$ and
$R=i$.  This honest signature and the forgery both pass all four challenge
classes and 80 large sampled challenges.  The same-message control matters:
it demonstrates that the primary forged message lies in a working honest
message domain and was not selected only because honest signing fails there.

\begin{table}[ht]
\centering
\caption{Primary full-parameter measurements.}
\begin{tabular}{@{}ll@{}}
\toprule
Item & Measurement \\
\midrule
Parameters & $q=1{,}048{,}609$, $d=2$, $r=3$, $m=448$, $s=13$ \\
Public coefficients & 2,240 \\
Signing queries & 0 \\
Public forger time & 0.001789767 s \\
Complete in-program time & 112.834013369 s \\
External wall time & 1:53.69 \\
Peak resident memory & 19,968 KiB \\
Workers & 1 \\
Honest same-message trials & 80/80 accepted \\
Forged fresh-message trials & 80/80 accepted \\
Exact challenge classes & 4/4 accepted \\
Negative controls & 7/7 rejected \\
\bottomrule
\end{tabular}
\end{table}

The pinned target and primary coefficient digests are:
\begin{quote}
\small
Target PDF:\\
\texttt{37ad500d8c0442173422d1976f9185b45302de819825624bdfa862de3ec5f9b9}\\
Public $F$:\\
\texttt{913281d14869095eb29d4a33c72e3c265364de7ee27ca4f0478f5fde998da160}\\
Forged $T^*$:\\
\texttt{656154385287f8445f17f1e184dd5e7d6b0d1eff5475afdcc9b20e6ea7c3234c}
\end{quote}

\subsection{Controls}

The seven negative controls remove the perturbation, replace it by each of
three wrong coordinate monomials, add two orbit-visible monomials to the valid
forgery, and reuse the forged signature on the changed message $E'=2j$.  Each
control either restores equality in the preliminary check or breaks the
challenge equality, and all seven reject.  Eight public-versus-secret
evaluation checks and eight checks of the public signed permutation also pass.

\subsection{Independent AI reproduction}

A separately written implementation reconstructed Equations~\eqref{eq:F} and
\eqref{eq:T} from the pinned target and used the different conforming prime
$q=1{,}048{,}583$.  It generated all 2,240 public coefficients, then invoked
the forger as an isolated process whose sole input was serialized public $F$.
That process had no signing-oracle or secret-material interface and recorded
zero queries.  The $g=2$ forgery passed all four challenge classes.  Only
after that acceptance, the implementation generated an honest signature on
the same $E=2k$; it also passed all four classes.  Omitting the perturbation,
using an orbit-visible perturbation, and changing the message to $E'=2j$ all
caused rejection.  The implementation also cross-checked both public
polynomials against separate direct evaluations of the paper's factored
equations.

The isolated public forger took 4.716746141 seconds, including JSON parsing
and serialization.  The complete run took 60.009294720 seconds internally
and 61.05 seconds externally, with 22,528~KiB peak resident memory and one
worker.  Its canonical coefficient digests are
\begin{quote}
\small
Public $F$:\\
\texttt{213ae3663e6cf2599d81c5141382bcc951699b8c6aee41712e10d66ce614441a}\\
Forged $T^*$:\\
\texttt{4a1350707a4bc3a043d5ab2d436b6214c826c3004994a06e2da08fab9b9cc633}\\
Honest same-message $T$:\\
\texttt{a00eefe7e020a573c1c11ee6ddaa9660f788b081bdea0f386c4f92c4e2db1d33}.
\end{quote}
The retained independent source, public-only forger, and 19-entry manifest
have SHA-256 digests, respectively,
\begin{center}
\small\ttfamily
4a82384994684fdc25f1240f0990f9aeb78a9826a629783f767fa31673f6580b\\
d6898781f9bb8a1e1a327978392a834e25e245e363755bdcea9da94d419cef03\\
2fb137a850f32d75334a0bcbcc848b1ffd59edf1390c8698cc622b601b917f7c.
\end{center}
The standalone public adaptation of the independent source has digest
\begin{center}
\small\ttfamily
27be188c8bd8042b43c243f9aa2f5977ab1b7266ad5ff49629e17bc6fe8ecb5a.
\end{center}

\section{Why the verifier fails}

The random challenge is much less random algebraically than its integer $p$
suggests.  With $R=i$, every challenge lies in a four-point set.  The
degree-13 monomial $x_2^{13}$ vanishes identically there.  The preliminary
check evaluates at $RE=-2j$, which is outside that set and has nonzero
$x_2$.  One coefficient therefore separates the two checks exactly.

Merely banning order-four $R$ is not a complete repair.  Every quaternion
$R$ satisfies a quadratic relation over its scalar field, so all powers of
$R$ lie in the at-most-two-dimensional commutative subalgebra
$\operatorname{span}_{\F_q}\{1,R\}$.  The challenge points consequently
remain in a predictable low-dimensional algebraic set even when $R$ has
large order.  The order-four choice gives the cleanest concrete full-row
forgery, but the structural issue is the restricted evaluation locus.

\section{Limitations and repair}

Our empirical claim is for the final revision and its explicit
$d=2,r=3,m=448$ row.  The proof is independent of the secret coefficients and
uses any odd 21-bit-or-larger prime, but we do not report measurements for
other rows.  The attack produces a direct forgery; it does not recover the
nominal secret key.  The paper gives no production implementation or complete
wire format, so the reproducer implements the mathematical coefficient
interface exactly as specified.  The key-sampling distribution is not fully
formalized; we both audit the hidden noncommutativity conditions on the test
key and give their probability under independent uniform sampling.  Finally,
the dated attack-history search cannot establish an absolute absence of
unindexed or private work.

A repair must replace the verification mechanism, rather than only blacklist
$i$.  In particular:
\begin{itemize}
  \item do not authenticate an attacker-supplied polynomial from evaluations
  confined to powers of an attacker-supplied quaternion;
  \item bind the signature to the message with a standard, analyzed signature
  transform and require evidence that the transmitted object was derived from
  the signing key;
  \item if polynomial identity testing remains part of verification, sample
  from a domain not contained in a predictable low-degree variety and prove a
  soundness bound for the actual polynomial class; and
  \item make all message-domain conditions publicly checkable and require
  honest signing correctness for every admitted message.
\end{itemize}
The single inequality~\eqref{eq:inequality} cannot provide these properties.

\section{Secondary correctness observation}

There is a simpler but weaker presentation of the same verifier flaw.  Taking
$E=k$ gives $g=1$, and
\[
 T^*(X)=F(X)+(x_2^{13},0,0,0)
\]
again passes every verifier challenge.  An honest signer on such a message,
however, computes $T=F$, which the preliminary inequality then rejects.
Thus the $g=1$ case also exposes a correctness defect in the stated message
domain.  We do not use it as the primary attack: the $g=2$ construction above
has an accepted honest signature and an accepted forgery on the same message.

\section{Conclusion}

For the claimed full parameter row, a zero-query adversary can turn the
public polynomial $F$ by one quaternion quarter turn and add one degree-13
coefficient.  The resulting $T^*$ passes Yagisawa's ordinary verifier for
every challenge while signing the fresh, honestly signable message $E=2k$.
The online attack takes about 1.8 milliseconds and requires no secret or
signature oracle.  The failure comes from testing an attacker-supplied
polynomial only on a structured power orbit and at one separately patched
point.

\section*{Acknowledgements}
The computations were performed on the Norwegian Research and Education Cloud (NREC), using resources provided by the University of Bergen and the University of Oslo.

\begin{thebibliography}{9}

\bibitem{yagisawa}
M.~Yagisawa.
\newblock A Digital Signature Using Multivariate Functions on Quaternion Ring.
\newblock IACR Cryptology ePrint Archive, Report 2010/352, final revision
\texttt{20100627:124304}, 2010.
\newblock \url{https://eprint.iacr.org/2010/352}.

\bibitem{zippel}
R.~Zippel.
\newblock Probabilistic algorithms for sparse polynomials.
\newblock In \emph{EUROSAM 1979}, Lecture Notes in Computer Science, vol.~72,
pp.~216--226. Springer, 1979.
\newblock \url{https://doi.org/10.1007/3-540-09519-5_73}.

\bibitem{schwartz}
J.~T.~Schwartz.
\newblock Fast probabilistic algorithms for verification of polynomial
identities.
\newblock \emph{Journal of the ACM}, 27(4):701--717, 1980.
\newblock \url{https://doi.org/10.1145/322217.322225}.

\end{thebibliography}

\end{document}
